\documentclass[11pt]{article}

\setlength{\topmargin}{-0.7in}
\setlength{\oddsidemargin}{-0.2in}
\setlength{\evensidemargin}{-0.2in}
\setlength{\textwidth}{6.9in}
\setlength{\textheight}{9.6in}
\setlength{\parindent}{0pt}
\setlength{\parskip}{0.6em}

\begin{document}

\begin{center}
    {\Large \textbf{Project Proposal - Data Management 2025/2026}}
\end{center}

\textbf{Group members}
\begin{itemize}
    \item Hamza Abdul Kader, matricola 2230150
    \item Leonardo Marasca, matricola 2217140
\end{itemize}

\textbf{Chosen type of project}

NoSQL / DBMS comparison: comparison between a graph database tool and a relational DBMS.

\textbf{Dataset}

OpenAlex dataset, accessed through the OpenAlex API. We will use a selected subset of AI-related research papers, authors, citations, topics, and abstracts.

Reference link: https://openalex.org/

Dataset documentation: https://help.openalex.org/hc/en-us/articles/24397285563671-About-the-data

\textbf{Project title}

Comparison between PostgreSQL and Neo4j for AI Research Paper Retrieval and Knowledge Graph Analysis

\textbf{Project description}

The goal of this project is to compare a relational DBMS and a graph database for the analysis of AI research papers. We intend to use a selected subset of the OpenAlex dataset focused on topics such as Artificial Intelligence, Machine Learning, and Retrieval-Augmented Generation. The data will include papers, authors, citations, topics, and abstracts.

We will model the same data in PostgreSQL using relational tables and in Neo4j using a graph structure. Then, we will perform analytical queries on both systems, such as finding highly cited papers, discovering connections between authors, analyzing citation paths, and retrieving papers related to specific AI topics. Finally, we will compare PostgreSQL and Neo4j in terms of data modeling, query complexity, expressiveness, and execution time.

\end{document}
